\scanpage{110}
% ms. line 01
P. ex. \deleted{\unsure{si}} $e^{(z^\omega)}$ est une fonction d'ordre $\omega$, de
% Paper watermark as on scan 109: MADE IN FRANCE (oval); The Best Typewriting Quality.
% ms. line 02
genre $\omega$. $\prod_n\left(1-\frac{z}{n^\alpha}\right)$ (où $\alpha>1$) est d'ordre
% ms. line 03
$1$ (sauf erreur) et de genre $0$.

% ms. line 04
Notons que l'ordre du produit de deux
% ms. line 05
fonctions entières est au plus égal au maximum
% ms. line 06
des deux ordres : $\omega(fg)\leq\mathrm{Max}(\omega f,\omega g)$, même
% ms. line 07
chose pour la somme. \deleted{D'où, \unsure{p. ex.} si $f(z)=e^{P(z)}\unsure{\prod}$}
% ms. line 07a
\inserted{l'ensemble des}
% ms. line 08
fonctions entières dont l'ordre est $<\mu$ donné
% ms. line 09
(resp\textsuperscript{t} $\leq\mu$) est une \deleted{\unsure{gr}} algèbre.

% ms. line 10
\uline{Théorème}3 (Hadamard) Soit la fonction entière
% ms. line 11
$f(z)=\sum c_n z^n$. Si $n^\alpha\sqrt[n]{|c_n|}\longrightarrow0$ pour
% ms. line 12
$n$ infini, alors $f$ est au plus de genre $\frac1\alpha$.
% archive sheet number: 111
